	\subsection{8. Conclusion}\label{conclusion}}

This paper presented A2C for proactive moving IoT service composition
with spatio-temporal constraints, preserving STR-based selection from
prior work. We compared shared and separate network architectures on the
same datasets (random waypoint and vehicle movement). A2C outperforms
Double DQN across all evaluated metrics: success rate (+19.2\% at 80
km/h), capacity satisfaction (+11.7\%), adaptation speed (2.3s vs 4.7s),
and re-composition frequency (69.2/hr vs 124.3/hr).

\textbf{Key Findings:}

\begin{enumerate}
	\def\labelenumi{\arabic{enumi}.}
	\item
	      \textbf{Proactive composition works}: Trajectory prediction enables
	      services to be selected that maintain adequate capacity throughout the
	      composition horizon, not just at the current moment.
	\item
	      \textbf{Architecture matters in high mobility}: Separate networks with
	      LSTM trajectory encoding achieve 73.5\% success rate at 80 km/h versus
	      54.3\% for Double DQN, a 35\% relative improvement.
	\item
	      \textbf{Shared networks offer efficiency}: For lower mobility
	      scenarios (2-5 km/h), shared networks converge faster (350 vs 500
	      episodes) while achieving 94-95\% success rate.
	\item
	      \textbf{Stability through entropy}: The entropy regularization
	      component reduces unnecessary re-compositions by 44\%, lowering system
	      overhead.
	\item
	      \textbf{Real data validates}: On the Illinois GPS dataset with 42,480
	      samples, the A2C agent achieves 92.3\% valid action selection,
	      demonstrating effective learning of capacity-based service selection.
\end{enumerate}

Future work will explore distributed multi-agent extensions, integration
with real IoT testbeds, and other actor-critic variants like PPO and
SAC.

\begin{center}\rule{0.5\linewidth}{0.5pt}\end{center}

\hypertarget{references}{%
